\scanpage{264}
\mspage{433}

\[
(ad-bc)D=(c^2+cd+d^2) \tag{23}
\]
thus $D$ \uline{is invertible \unsure{iff}} $c^2+\unsure{cd|d}+d^2=\mathcal{N}_{A_\rho/A}(c\rho+d)$ \uline{is i.e.} \unsure{iff} $c\rho+d$ is invertible, and then
\[
ad-bc=(c^2+cd+d^2)/D
\]
\margin{for $c^2+cd+d^2\equiv1\ (3)$}
which shows that the sign of $\det\alpha=ad-bc$ is then equal to the value of $D$ mod $3$.

(But can $D$ be equal to $0$ mod $3$? then mod $3$ the matrix $\xi$ \deleted{would be} would depend only on the parameters $A,B$
\[
\xi=\begin{pmatrix}A&B\\1+B&0\end{pmatrix}
\]
with $-B(1+B)=1$ \inserted{i.e. $B^2+B+1=0$} \deleted{i.e. $B=0$ or $B=-1$} \inserted{thus $B^3=1$ thus (since $3$) $B=1$}
thus we will have $\xi=\begin{pmatrix}A&1\\-1&0\end{pmatrix}$, \deleted{\illegible{4}}
\[
\xi\rho=\begin{pmatrix}-1&A+1\\0&-1\end{pmatrix},
\]
this matrix is regular \unsure{iff} $A\ne-1$, OK.\unsure{)|\}} Thus a priori we cannot exclude the case where $D$ \unsure{can} take a zero value mod $3$\ldots)

When we take the section $\sigma$ of $\unsure{\mathrm{Sl}(2)_{\mathbb{Z}}|\mathrm{Gl}(2)_{\mathbb{Z}}}$, given by
\[
\sigma=\sigma_\infty=\begin{pmatrix}0&-1\\1&0\end{pmatrix}
\qquad \sigma_\infty^4=1,\quad \sigma_\infty^2=-1 \tag{24}
\]
we again have (unless I am mistaken)
\margin{\unsure{case|for} $A=\mathbb{Z}$.}
\[
\mathrm{Pic}(A_\sigma)=\unsure{0|1}. \tag{25}
\]

\uline{The} numerical condition on
\[
\eta=\begin{pmatrix}R&S\\T&U\end{pmatrix} \tag{26}
\]
to \deleted{satisfy} be \unsure{locally} in the image

\scanpage{265}

of $F_\sigma$ (in addition to the \uline{regularity} condition on $\eta\sigma$)
\[
\eta=\beta\sigma(\beta)^{-1}=[\beta,\sigma] \tag{27}
\]
is that we have
\[
\det\eta=1,\quad \mathrm{Tr}(\eta\sigma)=\mathrm{Tr}\sigma \tag{28}
\]
now
\[
\eta\sigma=\begin{pmatrix}R&S\\T&U\end{pmatrix}
\begin{pmatrix}0&-1\\1&0\end{pmatrix}
=\begin{pmatrix}S&-R\\U&-T\end{pmatrix}
\]
and (28) becomes
\[
RU-ST=1,\quad S-T=0\quad\text{i.e. }S=T \tag{29}
\]
i.e. $\eta$ must be \uline{symmetric}. \deleted{\illegible{1}} If this condition is satisfied, it follows \unsure{from} (25) --- when the base is $\mathbb{Z}$ --- that we can find a solution $\beta\in\mathrm{Gl}(2,\mathbb{Z})$ of (27), the obstruction to finding a solution in $\mathrm{Sl}(2,\mathbb{Z})$ is again a sign, namely an element in
\[
\mathbb{Z}^*/\mathcal{N}(A_\sigma^*)\simeq\pm1
\]
--- as we have, for
\[
\beta=\begin{pmatrix}r&s\\t&u\end{pmatrix}. \tag{30}
\]
\deleted{\unsure{the relation}} satisfying (27), the relation\margin{\unsure{$\varepsilon$}}
\[
(\unsure{rt-su})U=t^2+u^2=\mathcal{N}(t\sigma+u) \tag{31}
\]

\scanpage{266}
\mspage{434}

we see as before that $U$ is invertible \unsure{iff}\margin{(arbitrary base)} $t\sigma+u$ is i.e. $t^2+u^2$ is, and that in this case, the obstruction \deleted{\unsure{of}} (base $\mathbb{Z}$, or $\mathbb{Z}_\ell$) is the residue of $U$ mod $4$ (which is equal to $\pm1$)\ldots
\margin{for $t^2+u^2\equiv1\ (4)$}

Let us summarize for reference

\uline{Proposition.} Let $A=\mathbb{Z}$, or $\mathbb{Z}_\ell$, or $\widehat{\mathbb{Z}}$, and let us consider the \inserted{two} equations \deleted{in}
\[
\left\{\begin{aligned}
\relax
[\alpha,\rho]&\overset{\unsure{\mathrm{déf}}}{=}\alpha\rho(\alpha)^{-1}=\xi\\
[\beta,\sigma]&=\beta\sigma(\beta)^{-1}=\eta
\end{aligned}\right.
\]
where $\rho=\begin{pmatrix}0&1\\-1&1\end{pmatrix}$, $\sigma=\begin{pmatrix}0&-1\\1&0\end{pmatrix}$,
$\xi$ and $\eta\in\mathrm{Gl}(2,A)$ given, with unknowns $\alpha$ resp. $\beta$ in $\mathrm{Gl}(2,A)$. We assume that
\[
\xi=\begin{pmatrix}A&B\\C&D\end{pmatrix}\quad\text{resp. }\eta=\begin{pmatrix}R&S\\T&U\end{pmatrix}
\]
satisfy the \inserted{necessary} numerical conditions for the existence of a solution
\[
AD-BC=1\qquad B+D-C=1 \tag{16}
\]
\[
RU-ST=1\qquad S-T=0\quad\text{i.e. }S=T. \tag{29}
\]

Then:

1) For there to be a solution \inserted{$\alpha$ (resp. $\beta$)} in $\mathrm{Gl}(2,A)$, it is \unsure{necessary|necessary} and sufficient that $\xi\rho=\begin{pmatrix}-B&A+B\\-D&C+D\end{pmatrix}$ be regular, i.e. nonscalar on each fiber (resp. that $\eta\sigma=\begin{pmatrix}S&-R\\U&-T=-S\end{pmatrix}$ be so) --- \deleted{and for}

\scanpage{267}

\deleted{\unsure{this it suffices}} i.e. that on no fiber, $\xi$ (resp. $\eta$) be a scalar multiple of $\rho^{-1}$ (resp. $\sigma^{-1}$) --- and for this it suffices that $D$ (resp. $U$) be invertible (since the \deleted{\unsure{coefficient}} \unsure{corresponding} \unsure{in|of}
\[
\rho^{-1}=\begin{pmatrix}1&-1\\1&0\end{pmatrix}\quad\text{or}\quad
\sigma^{-1}=\begin{pmatrix}0&1\\-1&0\end{pmatrix}
\]
is zero).

2) When $D$ (resp. $U$) is invertible, there is always a solution $\alpha$ resp. $\beta$, with $\det\alpha\in\pm1$ (resp. $\det\beta\in\{\pm1\}$).

3) For there to be a solution in $\mathrm{Sl}(2,A)$, it is necessary and sufficient that we have $D\equiv1\ (3)$ (resp. $U\equiv1\ (4)$) --- a condition which is automatic if $A=\mathbb{Z}_\ell$ with $\ell\ne3$ (resp. with $\ell\ne2$).

II) \uline{Study \inserted{in $\mathrm{Gl}(2)$ \deleted{\unsure{$\mathbb{Z}$}}} of the equation}
\[
[\alpha,\rho]=[\beta,\sigma]\unsure{\varepsilon_1}(\deleted{\illegible{1}}\mu-1). \tag{32}
\]
\margin{Study \unsure{first} of the equation in $\xi,\eta$.}

We always assume that $G$ is a form of $\mathrm{Gl}(2)$, associated with an Azumaya algebra $\mathcal M$ over the base $S=\mathrm{Spec}\,A$, and that we have \uline{two} \uline{regular} sections $\rho,\sigma$ of $G$ i.e. two invertible sections of $\mathcal M$. We are also given a ``$1$-parameter subgroup'' (additive) of $G$, i.e. a homomorphism
\[
\varepsilon_1:\mathbf G_a\longrightarrow \mathrm{S}G,\qquad \lambda\longmapsto\varepsilon_1(\lambda). \tag{33}
\]

\scanpage{268}
\mspage{435}

The case which interests us more particularly is that where $G=\mathrm{Gl}(2)_{\mathbb{Z}}$, $\rho$ and $\sigma$ as above, and
\[
\varepsilon_1(\lambda)=\varepsilon_1^\lambda=\begin{pmatrix}1&0\\\lambda&1\end{pmatrix}. \tag{34}
\]
We will note that here
\[
\varepsilon_1=\varepsilon_1(1)=\begin{pmatrix}1&0\\1&1\end{pmatrix}=\deleted{\unsure{\ell_1}}\rho\sigma \tag{35}
\]
We regard (32) as equations relating the quantities $\alpha,\beta,\mu$ --- $\alpha,\beta$ \inserted{variable local} sections of $G$, $\mu$ a section of $\unsure{\mathbb E^1|\mathbf E^1}$ --- the hypotheses we are going to make \unsure{imply|will imply} moreover that $\mu$ will in fact be a section of $\mathbf G_m$ (if, together with $\alpha,\beta$, it satisfies relation (32)). \deleted{\illegible{1}}

We will begin by studying the systems \deleted{of solutions} $(\xi,\eta,\unsure{\mu|\beta})$ satisfying
\[
\left\{\begin{aligned}
\xi&=\eta\varepsilon_1(\mu-1) && (33)\\
\det\xi&=\det\eta=1 && (34)\\
\mathrm{Tr}\xi\rho&=\underbrace{\mathrm{Tr}\rho}_{\inserted{\substack{\Vert\\\unsure{i}}}}\ \unsure{\text{and}}\quad
\mathrm{Tr}\eta\sigma=\underbrace{\mathrm{Tr}\sigma}_{\inserted{\substack{\Vert\\j}}} && (35)
\end{aligned}\right.
\]
\margin{\uline{N.B.} We will have $\det\varepsilon_1(\mu-1)=1$ thus $\det\xi=\det\eta$, \unsure{by the first relation}.}

We can eliminate $\eta$ from these equations, thanks to (33), we find

\scanpage{269}

\[
(36)\qquad
\det\xi=1,\qquad
\operatorname{Tr}(\xi\inserted{-1})\rho
=\deleted{\unsure{i}}\inserted{0},\qquad
\unsure{\operatorname{Tr}(\xi\varepsilon_1(1-\mu)\deleted{\sigma}\inserted{-1})\inserted{\sigma}}
\deleted{=\unsure{j}}=0.
\]

\deleted{where we have \unsure{set}, to \unsure{simplify}
\[
i=\operatorname{Tr}\rho\qquad j=\operatorname{Tr}\sigma
\]}

The third relation (36), on \unsure{expanding}
\unsure{with respect to the} coefficients of $\xi$, \deleted{and} can be written
\[
(37)\qquad P(\xi)\mu\unsure{-|+}Q(\xi)=0
\]
where the coefficients $P,Q$ are affine linear functions with respect to those of
$\xi$, \deleted{and} \unsure{give} rise to the identity (in $\xi,\mu$)
\[
(38)\qquad
\unsure{\bigl(\operatorname{Tr}(\xi\varepsilon_1(1-\mu)\deleted{\sigma}\inserted{-1})\inserted{\sigma}\deleted{-j}\bigr)}
=P(\xi)\mu-Q(\xi)\deleted{\unsure{=0}}
\]
\deleted{\unsure{$\operatorname{Tr}(\sigma\xi\varepsilon_1(1-\mu))$}}

The constant term $-Q(\xi)$ being \deleted{\unsure{equal to}}
\inserted{identically} \inserted{equal to} the value of the left-hand
\unsure{side} of (38) for $\mu=0$, i.e.
\[
\begin{aligned}
(39)\qquad -Q(\xi)
&=\operatorname{Tr}((\xi\varepsilon_1-1)\sigma).\\
&=\operatorname{Tr}\xi(\varepsilon_1\sigma)-\operatorname{Tr}\sigma
\end{aligned}
\]

We assume that we have
\[
(40)\qquad \varepsilon_1\sigma=-\lambda_0\rho
\qquad
\begin{gathered}
\lambda\text{ central section of }G\\
\text{i.e. }\lambda\text{ invertible scalar}
\end{gathered}
\]
so that
\[
(40\text{ bis})\qquad \boxed{\varepsilon_1=-\lambda_0\rho\sigma^{-1}}
\]
\margin{\uline{N.B.} If $\det\rho=\det\sigma=1$,
then we must have $\lambda_0^2=1$
thus if $A$ \unsure{integral}, $\lambda_0\in\{\pm1\}$.}
and (39) becomes, by virtue \deleted{of \unsure{the equation}}
\[
\deleted{\unsure{-}Q(\xi)=}\qquad
\operatorname{Tr}\xi(\varepsilon_1\sigma)
=\unsure{-\lambda_0|\lambda_0}\operatorname{Tr}\xi\rho
\overset{\inserted{\text{by }(36)}}{=}
-\lambda_0\operatorname{Tr}\rho,
\quad\text{thus}
\]

\scanpage{270}
\mspage{436}

we find
\[
\deleted{(\unsure{42|37})}\ (41)\qquad
Q(\xi)=\underbrace{\lambda_0\operatorname{Tr}\rho+\operatorname{Tr}\sigma}_{c}
\qquad
(\text{constant, independent of }\xi)
\]
and we assume
\[
(42)\qquad
c=\lambda_0\operatorname{Tr}\rho+\operatorname{Tr}\sigma
\quad
\deleted{\text{\unsure{integer|section} \unsure{element}}}
\inserted{\text{section}}\text{ \uline{invertible} of }\mathcal O_S.
\]

(In the case which interests us, we will have
\[
\sigma^{-1}=-\sigma,\quad
\varepsilon_1=\rho\sigma=-\rho\sigma^{-1},\quad
\text{i.e. }\lambda_0=1
\]
\[
\deleted{\unsure{\operatorname{Tr}}}\quad
c=\lambda_0\underbrace{\operatorname{Tr}\rho}_{1}
+\underbrace{\operatorname{Tr}\sigma}_{0}=1
\qquad\text{OK.}
\]
)

This being so, the \inserted{\uline{third}} equation (36) becomes
\unsure{(taking account of the first two)}
\[
(43)\qquad P(\xi)\mu=c
\]
which also means that $P(\xi)$ \uline{is invertible}, and
\[
(44)\qquad \mu=c/P(\xi).
\]

Thus, $\mu$ is uniquely determined by $\xi$, by equation (44) ---
and $\mu$ will necessarily be invertible.

In summary:

\uline{Proposition} If we assume that we have \deleted{(40 bis)}
an expression (40 bis) for the unipotent element
$\varepsilon_1=\varepsilon_1(1)$, and that we have (42)
($\lambda_0\operatorname{Tr}\rho+\operatorname{Tr}\sigma$ invertible), then

\scanpage{271}

the set of solutions \deleted{\unsure{of the|in}} $(\xi,\eta,\mu)$ of the
simultaneous equations (33)(34)(35)
is in $1$-$1$ correspondence with the set of $\xi$
\unsure{sections|section} of $G=G(\mathcal{M})$
satisfying the conditions
\[
(45)\qquad
\left\{
\begin{gathered}
\det\xi=1,\quad \operatorname{Tr}\xi\rho=\operatorname{Tr}\rho,\quad\text{and}\\
P(\xi)\text{ invertible section of }\mathcal O_S.
\end{gathered}
\right.
\]

In terms of $\xi$, the corresponding solution of the \unsure{system} (33 : 35)
is given by
\[
(\xi,\ \eta=\xi\varepsilon_1(1-\mu),\ \mu=c/P(\xi))
\]
--- and $\mu$ (the ``multiplier'') is necessarily invertible.

\uline{N.B.} If we had solved for $\eta$, we would have
found a condition of the form
\[
P'(\eta)\mu-Q'(\eta)=0
\]
with
\[
Q'(\eta)=\operatorname{Tr}\eta\varepsilon_0^{-1}\rho-\operatorname{Tr}\rho
\]
and the natural condition would be
\[
\varepsilon_0^{-1}\rho=-\lambda'_0\sigma
\qquad(\lambda'_0\text{ invertible scalar}),
\]
which is equivalent to (40 bis) with $\lambda'_0=1/\lambda_0$,
and implies that
\[
Q'(\eta)=c^{\mathrm{te}}
=\lambda'_0\operatorname{Tr}\sigma-\operatorname{Tr}\rho
=\deleted{\unsure{\lambda'_0}}c',
\]

\scanpage{272}
\mspage{437}

we see at once that
\[
c'=\lambda'_0c
\]
thus the condition $c$ invertible is equivalent to $c'$ invertible,
and we will then have the system of equations in $\eta$
\[
\det\eta=1,\quad
\operatorname{Tr}\eta\sigma=\operatorname{Tr}\sigma,\quad
P'(\eta)\mu=c'\quad\text{i.e. }
\left\{
\begin{gathered}
P'(\eta)\text{ invertible}\\
\mu=c'/P'(\eta)
\end{gathered}
\right.
.
\]

We see therefore that the hypotheses made are in fact symmetric in
$\xi,\eta$, and allow a symmetric treatment of the solution,
either via $\xi$, or via $\eta$.

Determination of $P(\xi)$. We will have
\[
\varepsilon_1(\lambda)=1+N_1\lambda=(1+N_1)^\lambda
\]
$N_1$ a square-zero matrix
\[
N_1^2=1,
\]
and
\[
\operatorname{Tr}(\xi\varepsilon_1(1-\mu)\sigma)-\operatorname{Tr}\sigma
=\operatorname{Tr}\xi\sigma-\operatorname{Tr}\sigma
+\unsure{\lambda}(1-\mu)\operatorname{Tr}\xi N_1\sigma
\]
thus
\[
(47)\qquad
P(\xi)=\deleted{\unsure{-}}\operatorname{Tr}\xi
\underbrace{(-N_1\sigma)}_{\inserted{\substack{\Vert\\N_1\rho=\rho+\sigma}}}
\]
Moreover \deleted{the formula} we will have
\[
(48)\qquad \varepsilon_1=1+N_1
\]
and formula (40 bis) reads
\[
\deleted{\sigma+N_1\sigma=-\lambda_0\rho}
\]
\[
(49)\qquad N_1=-(\lambda_0\rho\sigma^{-1}+1)
\]
\margin{\unsure{i.e.}}
\[
\deleted{N_1\sigma=-\lambda_0\rho-\sigma}
\]
\deleted{i.e.} \unsure{or} again

\scanpage{273}

\[
N_1\sigma=-(\lambda_0\rho+\sigma)
\]
thus we will have
\[
(50)\qquad
P(\xi)=\operatorname{Tr}\deleted{\unsure{\xi}}\xi(\lambda_0\rho+\sigma)
\]

\uline{Corollary to the Proposition} The scheme of simultaneous solutions
\inserted{of the equations} (33) : (35)
\inserted{in $\xi,\eta,\mu$} is a smooth relative scheme,
of relative dimension $2$ --- for the equation $\operatorname{Tr}(\xi\rho)=c^{\mathrm{te}}$
always defines a smooth subscheme of dimension $1$
of \deleted{\unsure{$\mathrm{Sl}(2)$|SG}} $SG(\mathcal{M})$
(as soon as $\rho$ is a section that is zero on no fiber),
--- of which we must then take the open subset defined by
``$P(\xi)$ invertible''.

In the case which interests us,
\[
\rho=\begin{pmatrix}0&1\\-1&1\end{pmatrix},\qquad
\sigma=\begin{pmatrix}0&-1\\1&0\end{pmatrix},
\]
\deleted{the condition in $\xi$ \illegible{1} $P(\xi)$ invertible reads}

\unsure{setting|for}
\[
\xi=\begin{pmatrix}A&B\\C&D\end{pmatrix},\qquad
\eta=\begin{pmatrix}R&S\\T&U\end{pmatrix}
\]
we will have
\[
P(\xi)=D,\qquad P'(\eta)=U
\]
\[
(51)\qquad
\unsure{\eta}\varepsilon_1(\mu-1)
=\begin{pmatrix}
R+(\mu-1)S,&S\\
T+(\mu-1)U&U
\end{pmatrix}
\]